\chapter{Valuations, valuation rings, and places}
\label{par-chap-foundations}
\source{papaioannou-algebraic-rr:page-0002}\uses{}
Throughout, $\mathbb K$ is a field. A function field is understood in the
one-variable sense used by the source.

\begin{definition}[Function field of one variable]
\label{par-def-function-field}
\dcref{1.1}
\source{papaioannou-algebraic-rr:page-0002}\uses{}
A function field $\mathbb F/\mathbb K$ of one variable is an algebraic extension
$\mathbb F\supseteq\mathbb K$ such that $\mathbb F$ is a finite extension of
$\mathbb K(x)$ for an indeterminate $x$ (equivalently, for an element of
$\mathbb F$ transcendental over $\mathbb K$). The algebraic closure $\widetilde{\mathbb K}$
of $\mathbb K$ in $\mathbb F$ is the field of constants. It is itself a function
field over $\mathbb K$; when $\widetilde{\mathbb K}=\mathbb K$, $\mathbb K$ is
the full constant field of $\mathbb F$.
\end{definition}

\begin{definition}[Valuation ring]
\label{par-def-valuation-ring}
\dcref{1.2}
\source{papaioannou-algebraic-rr:page-0002}\uses{}
A valuation ring of $\mathbb F/\mathbb K$ is a ring $\mathbb K\subsetneq\mathcal O
\subsetneq\mathbb F$ such that $z\notin\mathcal O$ implies $z^{-1}\in\mathcal O$.
\end{definition}

\begin{theorem}[Basic properties of a valuation ring]
\label{par-thm-valuation-ring-basic}
\dcref{1.1}
\uses{par-def-function-field,par-def-valuation-ring}
\source{papaioannou-algebraic-rr:page-0002}\uses{}
Let $\mathcal O$ be a valuation ring of $\mathbb F/\mathbb K$. Then:
\begin{enumerate}
\item $\mathcal O$ is local, with unique maximal ideal $\mathfrak m$ equal to
the complement of $\mathcal O^\times$;
\item if $\widetilde{\mathbb K}$ is the field of constants, then
$\widetilde{\mathbb K}\subseteq\mathcal O$ and
$\widetilde{\mathbb K}\cap\mathfrak m=\{0\}$.
\end{enumerate}
\end{theorem}
\begin{proof}
Put $\mathfrak m=\mathcal O-\mathcal O^\times$. It is closed under
multiplication by elements of $\mathcal O$. If $x,y\in\mathfrak m$ are nonzero,
one of $x/y,y/x$ lies in $\mathcal O$; say $x/y\in\mathcal O$. Then
$x+y=y(1+x/y)\in\mathcal O$, proving that $\mathfrak m$ is an ideal and hence
the unique maximal ideal. If $z\in\widetilde{\mathbb K}$ were outside $\mathcal O$,
then $z^{-1}\in\mathcal O$ and a relation
$a_mz^{-m}+\cdots+a_1z^{-1}+1=0$ over $\mathbb K$ would give
$z=-(a_mz^{-m+1}+\cdots+a_1)\in\mathcal O$, a contradiction. Since every
nonzero constant is a unit of $\mathcal O$, its intersection with the maximal
ideal is $\{0\}$.
\end{proof}

\begin{lemma}[Descending chains in the maximal ideal]
\label{par-lem-maximal-chain}
\dcref{1.3}
\uses{par-thm-valuation-ring-basic}
\source{papaioannou-algebraic-rr:page-0002}\uses{}
If $x\in\mathfrak m$ and $x_1=x,x_2,\ldots,x_n\in\mathfrak m$ satisfy
$x_i\in x_{i+1}\mathfrak m$, then
$n\leq[\mathbb F:\mathbb K(x)]<\infty$.
\end{lemma}
\begin{proof}
If a nontrivial relation $\sum a_ix_i=0$ over $\mathbb K(x)$ existed, clear
denominators and take relatively prime polynomial coefficients. Let $b_i=a_i(0)$
and choose maximal $j$ with $b_j\neq0$ and $b_i=0$ for $i>j$. Rewriting the
relation and dividing by $x_j$ shows $-a_j\in\mathfrak m$, while
$a_j=b_j+xg_j$ gives $b_j\in\mathfrak m\cap\mathbb K$, contradicting the
second part of Theorem~1.1. Thus the $x_i$ are $\mathbb K(x)$-independent, and
the finite degree bound follows from the function-field definition.
\end{proof}

\begin{theorem}[Structure of a valuation ring]
\label{par-thm-valuation-ring-structure}
\dcref{1.2}
\uses{par-def-valuation-ring,par-lem-maximal-chain}
\source{papaioannou-algebraic-rr:page-0002}\uses{}
If $\mathcal O$ is a valuation ring of $\mathbb F/\mathbb K$ with maximal
ideal $\mathfrak m$, then:
\begin{enumerate}
\item $\mathfrak m$ is principal;
\item if $t$ generates $\mathfrak m$, every $x\in\mathbb F$ has a representation
$x=ut^n$ with $u\in\mathcal O^\times$;
\item $\mathcal O$ is a principal ideal domain, and every ideal $\mathfrak n$ is
$t^n\mathcal O$ for some integer $n$.
\end{enumerate}
\end{theorem}
\begin{proof}
If $\mathfrak m$ were not principal, repeatedly choose
$x_{i+1}\in\mathfrak m-x_i\mathcal O$ starting with $x_1\in\mathfrak m$.
The valuation-ring alternative gives $x_i\in x_{i+1}\mathfrak m$, contradicting
the preceding lemma. For (b), after replacing $x$ by $x^{-1}$ assume $x\in\mathcal O$.
If $x$ is a unit take $n=0$; otherwise choose the largest $n$ with
$x\in t^n\mathcal O$, obtaining $x=ut^n$ with $u$ a unit. For (c), if
$\mathfrak n\neq0$, choose the minimum exponent among representations of its
nonzero elements; then $\mathfrak n=t^n\mathcal O$.
\end{proof}

\begin{definition}[Place]
\label{par-def-place}
\dcref{1.3}
\source{papaioannou-algebraic-rr:page-0003}\uses{}
A place $P$ of $\mathbb F/\mathbb K$ is the maximal ideal of a valuation ring
$\mathcal O$ of $\mathbb F/\mathbb K$. A generator $t$ of $P$ is a parameter.
The set of all places is denoted $\mathbf P_{\mathbb F}$.
\end{definition}

\begin{definition}[Discrete valuation]
\label{par-def-valuation}
\dcref{1.4}
\source{papaioannou-algebraic-rr:page-0003}\uses{}
A valuation is a surjective homomorphism
$v:\mathbb F^*\twoheadrightarrow\mathbb Z$ such that
$v|_{\mathbb K^*}=0$ and $v(x+y)\geq\min\{v(x),v(y)\}$. Set $v(0)=\infty$.
\end{definition}

\begin{theorem}[Strict triangle inequality when values differ]
\label{par-thm-valuation-strict}
\dcref{1.4}
\uses{par-def-valuation}
\source{papaioannou-algebraic-rr:page-0004}\uses{}
If $v(x)\neq v(y)$, then $v(x+y)=\min\{v(x),v(y)\}$.
\end{theorem}
\begin{proof}
Assume $v(x)<v(y)$. If $v(x+y)\neq v(x)$, then applying the triangle inequality
to $(x+y)-y$ gives
$v(x)=v((x+y)-y)>\min\{v(x+y),v(y)\}>v(x)$, a contradiction.
\end{proof}

\begin{definition}[Place valuation and residue field]
\label{par-def-place-valuation}
\uses{par-def-place,par-def-valuation,par-thm-valuation-ring-structure}
\source{papaioannou-algebraic-rr:page-0004}\uses{}
For $P$ with parameter $t$, write each nonzero $x\in\mathbb F$ uniquely as
$x=ut^m$ with $u$ a unit and set $v_P(x)=m$, $v_P(0)=\infty$. This is
independent of $t$. Conversely $v$ gives
$P=\{z:v(z)>0\}$ and $\mathcal O_P=\{z:v(z)\geq0\}$, with exact sequence
\[
0\longrightarrow\mathcal O_P^\times\longrightarrow\mathbb F^*
\xrightarrow{v}\mathbb Z\longrightarrow0.
\]
The valuation rings, valuations, and places are therefore equivalent discrete
categories. The residue map $\pi:\mathcal O_P\to\mathcal O_P/P$ extends by
the value $\infty$ outside $\mathcal O_P$; $\mathbb K$ embeds in the residue
field $\mathbb F_P:=\mathcal O_P/P$. The degree is $\deg P=[\mathbb F_P:\mathbb K]$.
\end{definition}

\begin{theorem}[Degree bound for a place]
\label{par-thm-place-degree}
\dcref{1.5}
\uses{par-def-place,par-def-place-valuation,par-def-function-field}
\source{papaioannou-algebraic-rr:page-0004}\uses{}
If $P$ is a place and $x\in P^*$, then
\[
\deg P\leq[\mathbb F:\mathbb K(x)]<\infty.
\]
\end{theorem}
\begin{proof}
If residue classes $z_i(P)$ of elements $z_i\in\mathcal O_P$ are
$\mathbb K$-independent but the $z_i$ are dependent over $\mathbb K(x)$,
clear denominators and write coefficients $a_i=b_i+xg_i$ with not all $b_i=0$.
Reducing modulo $P$ gives $\sum b_i z_i(P)=0$, a contradiction. The finiteness
is the finite-extension clause in the function-field definition.
\end{proof}

\begin{definition}[Zeros and poles]
\label{par-def-zero-pole}
\dcref{1.5}
\uses{par-def-place-valuation}
\source{papaioannou-algebraic-rr:page-0005}\uses{}
For $z\in\mathbb F^*$ and $P\in\mathbf P_{\mathbb F}$, $P$ is a zero of order
$m>0$ when $v_P(z)=m$, and a pole of order $-m$ when $v_P(z)=m<0$.
\end{definition}

\begin{theorem}[Existence of a place containing an ideal]
\label{par-thm-place-existence}
\dcref{1.6}
\uses{par-def-function-field,par-def-place,par-def-valuation-ring}
\source{papaioannou-algebraic-rr:page-0005}\uses{}
Let $R$ be a subring with $\mathbb K\subseteq R\subseteq\mathbb F$ and let
$\mathfrak a$ be a proper ideal of $R$. There is a place $P$ with
$\mathfrak a\subseteq P$ and $R\subseteq\mathcal O_P$.
\end{theorem}
\begin{proof}
Order by inclusion the set of subrings $S$ containing $R$ for which
$\mathfrak aS\neq S$. Zorn's lemma gives a maximal such $\mathcal O$.
If neither $z$ nor $z^{-1}$ belonged to $\mathcal O$, maximality would give
relations $1=\sum a_i z^i$ and $1=\sum b_jz^{-j}$ with coefficients in
$\mathfrak a\mathcal O$. Choosing minimal exponents and combining the relations
produces a relation of smaller exponent, contradiction. Hence $\mathcal O$ is a
valuation ring and its maximal ideal contains $\mathfrak a$.
\end{proof}

\begin{corollary}[Every transcendental element has a zero and a pole]
\label{par-cor-zero-pole}
\dcref{1.7}
\uses{par-thm-place-existence,par-def-zero-pole}
\source{papaioannou-algebraic-rr:page-0005}\uses{}
If $z$ is transcendental over $\mathbb K$, then $z$ has a zero and a pole, so
$\mathbf P_{\mathbb F}\neq\varnothing$.
\end{corollary}
\begin{proof}
Apply the theorem to $R=\mathbb K[z]$ and $\mathfrak a=z\mathbb K[z]$ to obtain
a zero of $z$; apply it to $z^{-1}$ to obtain a zero of $z^{-1}$, hence a pole of $z$.
\end{proof}

\begin{corollary}[Finite constant-field extension]
\label{par-cor-constant-field-finite}
\dcref{1.8}
\uses{par-cor-zero-pole,par-thm-place-degree}
\source{papaioannou-algebraic-rr:page-0005}\uses{}
The field of constants $\widetilde{\mathbb K}$ is a finite algebraic extension of
$\mathbb K$.
\end{corollary}
\begin{proof}
Choose a place. The residue embedding of $\widetilde{\mathbb K}$ into
$\mathbb F_P$ gives
$[\widetilde{\mathbb K}:\mathbb K]\leq[\mathbb F_P:\mathbb K]<\infty$.
\end{proof}
